\documentclass[12pt]{article}

\usepackage[a4paper, total={6in, 8in}]{geometry}
\usepackage{textcomp}

\begin{document}
\setlength{\parindent}{0pt}

\def \t {\textrightarrow}
\def \v {\vspace{1em}}
\def \bi {\begin{itemize}}
\def \ei {\end{itemize}}
\def \s[#1] {\section*{#1}}
\def \ss[#1] {\subsection*{#1}}
\def \sss[#1] {\subsubsection*{#1}}

\s[D'annunzio]
\ss[Il fuoco]
Il fuoco \t passione correlata al V canto dell'inferno \\
La passione diventa per lui una vera nevrosi \t voleva riprodurre un'opera d'arte che inglobasse tutti i campi \\
Questo genera la fase di ripiegamento (depressione), la delusione e l'inettitudine \t questo lo vede lontano da Nietzsche \\
Non c'è chiusura del cerchio \t frana tutto, e il romanzo di d'annunzio no si ferma all'estetismo bruto \t ma getta i semi per i romanzi del 900 \\
Questo è stato visto come una simulazione del fuoco in pittura \t Alciati, omicidio-suicidio di due che abbriacciati si buttano \\
Nel fuoco la delusione si fa fortissima \t per ragionare di un d'annunzio che non è solo il piacere, e che il Vittoriale e La capponcina descrivono \\
C'è anche il momento in cui si ritrova con la delusione del proprio animo \t Ezio Raimondi dice, 1969 "Una vita come opera d'arte" che d'annunzio sarebbe stato gradito a un pubblico moderno, e parla anche delle volgarita dell'autore \\
L'essere volgari comprende come ci si espone al pubblico \\
Nella sua visione di inettitudine che trapela nella situazione presentata nel fuoco \\
Anche nel fuoco c'è una compresenza di aspetti \t dalla parte biografica alla parte dell'opera \\
Rottura con la Duse che è stata anticipata in questo romanzo

\v

Un opera in versi, romanzo in versi \\
La Gravina e la Leoni \t sono donne che d'annunzio ha avuto \\
Nel 97 entra in politica \t volteggia tra destra e sinistra, e paradigma con Mussolini \t ricordare:
\bi
  \item il notturno
  \item poema paridisiaco
  \item il diario segreto
\ei
Il suo voyerismo lo porta a sfidare il vuoto, vola su vienna e perde la vista da un occhio \t questo ha menomato la sua capacita visiva (ma non quella cognitiva) \t scrive pero con uno spirito diverso: fine della sua parabola \\
Nel 31 si ritira sul garda \t morira nel 38 \\
I rapporti dal 22 con il regime fascita sono ambivalenti: da un lato di unione, dall'altro delusione per il trattamento che riceve da mussolini \\
D'annunzio non approva pero l'alleanza con la germania \t e definisce Hitler "il ridicolo nibelungo truccato alla Charleu" \t la rottura si verifica con l'avvicinamento di Mussolini alla germania \\

\ss[Forse che si forse che no]
Sembra quasi il nome di un opera di Pirandello \t il dubbio pirandelliano \\
Il protagonista di quest opera vuole sperimentare il nuovo secolo attraverso gli idoli di questo secolo, ovvero la velocita, la macchina e il veivolo \\
Arditismo e andare oltre i propri limiti sono presenti \t paradigma della volonta di andare oltre il progresso \\
Il tema del progresso che è stato enunciato quando abbiamo paralato di leopardi \t e anche con verga che si trova restio al progresso, e Compte afferma che senza la scienza non c'è progresso \\
Verga fa seguito a Compte, pero arriva a una negativa considerazione del progresso \t l'ostrica via dallo scoglio trova la morte \\
Un progresso che avvicina ma allontana \t Carducci in "La stazione di un mattino d'autunno" \t un marchingegno che avvicina (permette il trasporto veloce), ma allontana perche le perosne possono spostarsi \t l'amata dall'autore parte con il treno \\
Il romanzo che inquadra il progresso e il salto nel vuoto con il volo sopra vienna

\v

La fase notturna è una fase di ripiegamento interiore \t anche se l'attivita sessuale rimane attiva \\
La nevrosi di matrice sessuale lo porta ad avere un forte narcisismo \\
Siamo gia nella fase del declino di d'annunzio \t e la nuova delusione ha un nome nuovo, che rallegra la sua anima: Luisa Baccara, è una pianista, che condivide con lui dei momenti di piacevole svago al Vittoriale \t a Gardone Riviera \\
Muore il 18 marzo del 38 da un emorragia celebrale \\
Il fenomeno d'annunzio fa scena al tempo \t ci si aspetta questo personaggio, in accordo con la filosofia di Nietzsche \t quasi propagandista \\
È contrapposto a Pascoli \t il nido, affetti semplici e morbosita \t D'annunzio è voyerismo etc. \\
Siamo anche a ridosso della prima guerra mondiale \t si dibatte sull'interventismo o il neutralismo \t partecipare in guerra significava un certo coraggio \\
Simmel \t la moda del tempo
\end{document}
